\subsection{Composition of bounded diffusions}

PROPOSITION 13. --- Let T,X,Y be three locally compact spaces, \(\Lambda : t \mapsto \lambda_t\) a bounded diffusion of T in X, and H: \(x \mapsto \eta_x\) a bounded diffusion of X in Y. The mapping \(t \mapsto \lambda_t\mathrm{H}\) is then a bounded diffusion of T in Y, which is denoted by \(\Lambda\mathrm{H}\) , and

\[
\| \Lambda H \| \leqslant \| \Lambda \| \| H \|.\tag{14}
\]

Let \(f\) be a universally measurable function \(\geqslant 0\) defined on \(Y\) , and \(\mu\) a measure on \(T\) . Suppose that \(\mu\) belongs to the domain of \(\Lambda\) , and that \(\mu \Lambda\) belongs to the domain of H; then \(\mu\) belongs to the domain of \(\Lambda \mathrm{H}\) , and

\[
\begin{array}{l} \langle \mu (\Lambda \mathrm{H}), f \rangle = \langle \mu \Lambda , \mathrm{H} f \rangle = \langle \mu , \Lambda \mathrm{H} f \rangle ; \\ (\mu \Lambda) \mathrm{H} = \mu (\Lambda \mathrm{H}); \quad \Lambda (\mathrm{H} f) = (\Lambda \mathrm{H}) f. \end{array}\tag{15}
\]

Set \(\gamma_{t} = \lambda_{t}H\) ; we shall denote by \(\Gamma\) the mapping \(\Lambda H\) of T into \(\mathcal{M}_{+}(Y)\) , and by \(\Gamma f\) the function \(t \mapsto \langle \gamma_{t}, f \rangle\) (an abuse of notation, since we do not yet know whether \(\Gamma\) is a diffusion). Then \(\langle \gamma_{t}, f \rangle = \langle \lambda_{t}H, f \rangle = \langle \lambda_{t}, Hf \rangle\) by (13); since the function Hf is positive and universally measurable on X (Cor. of Prop. 12), it follows first of all that \(\Gamma f = \Lambda(Hf)\) , and then that \(\Gamma f\) is universally measurable on T (same reference). It is clear that all of the measures \(\gamma_{t}\) have total mass at most equal to \(\|\Lambda\| \|H\|\) . Consequently \(\Gamma g\) is universally measurable and bounded for every function \(g \in \mathcal{K}_{+}(Y)\) ; \(\Gamma\) is therefore scalarly essentially integrable for every bounded measure on T, and in particular for every measure with compact support. More generally, if \(\mu\) is a measure in the domain of \(\Lambda\) , such that \(\mu\Lambda\) belongs to the domain of H, then, for \(g \in \mathcal{K}_{+}(Y)\) ,

\[
\langle \mu , \Gamma g \rangle = \langle \mu , \Lambda (\mathrm{H} g) \rangle = \langle \mu \Lambda , \mathrm{H} g \rangle = \langle (\mu \Lambda) \mathrm{H}, g \rangle .
\]

Since the last quantity is finite, we see that \(\Gamma\) is scalarly essentially \(\mu\) -integrable. Let us denote by \(\mu\Gamma\) the integral \(\int\gamma_{t}d\mu(t)\) (an abuse of notation, since we do not yet know whether \(\Gamma\) is a diffusion). The preceding relations may then be written

\[
\langle \mu \Gamma , g \rangle = \langle (\mu \Lambda) \mathrm{H}, g \rangle ,
\]

or also \(\mu \Gamma = (\mu \Lambda)\mathrm{H}\) since \(g\) is arbitrary in \(\mathcal{K}_+(\mathrm{Y})\) .

Consider anew the universally measurable function \(f \geqslant 0\) . We have

\[
\langle \mu \Gamma , f \rangle = \langle (\mu \Lambda) \mathrm{H}, f \rangle = \langle \mu \Lambda , \mathrm{H} f \rangle = \langle \mu , \Lambda (\mathrm{H} f) \rangle = \langle \mu , \Gamma f \rangle .
\]

When f is lower semi-continuous and \(\mu\) runs over the set of positive measures with compact support, these relations express that \(\Gamma\) is a diffusion of T in Y. The assertion then does no more than make explicit the relations obtained in the course of the above proof.

DEFINITION 4. --- The notations being those of Proposition 13, the diffusion \(\Lambda H\) is called the composed diffusion (or the composition) of the bounded diffusions \(H\) and \(\Lambda\) .

Let \(X_{1}, X_{2}, X_{3}, X_{4}\) be four locally compact spaces, and \(\Lambda_{1}, \Lambda_{2}, \Lambda_{3}\) three bounded diffusions, of \(X_{1}\) in \(X_{2}\) , \(X_{2}\) in \(X_{3}\) , \(X_{3}\) in \(X_{4}\) , respectively. It follows at once from Prop. 13 that

\[
(\Lambda_ {1} \Lambda_ {2}) \Lambda_ {3} = \Lambda_ {1} (\Lambda_ {2} \Lambda_ {3}).
\]

We will therefore use these notations without parentheses for the composition of diffusions.

Example. --- Let u be a universally measurable mapping of T into X, and v a universally measurable mapping of X into Y; by Prop. 2 b), one defines diffusions \(\Lambda\) and H by the formulas

\[
\lambda_ {t} = \varepsilon_ {u (t)}, \eta_ {x} = \varepsilon_ {v (x)};
\]

the diffusion \(\Gamma = \Lambda H\) is then given by

\[
\gamma_ {t} = \varepsilon_ {(v \circ u) (t)}.
\]

One is therefore careful to note that the order of composition of diffusions is the opposite of the usual order of the composition of functions.

